\begin{table}[t]
\centering
\small
\begin{tabular}{lrrrr}
\toprule
& \multicolumn{3}{c}{\textbf{Replacement}} & \textbf{Augment.} \\
\cmidrule(lr){2-4}\cmidrule(lr){5-5}
\textbf{Input representation} & \textbf{Macro-F1} & \textbf{Wtd-F1} & \textbf{Acc.} & \textbf{Macro-F1} \\
\midrule
\multicolumn{5}{l}{\textit{Tamil}} \\
\quad No transliteration & 52.91 & 56.00 & 51.78 & 52.91 \\
\quad Baseline transliteration & 52.69 & 54.98 & 50.59 & 52.89 \\
\quad Phonology-aware transliteration & 52.22 & 55.37 & 50.99 & 52.43 \\
\addlinespace
\multicolumn{5}{l}{\textit{Malayalam}} \\
\quad No transliteration & 69.86 & 68.60 & 67.84 & 69.86 \\
\quad Baseline transliteration & 71.17 & 69.60 & 68.95 & 70.23 \\
\quad Phonology-aware transliteration & 70.99 & 69.31 & 68.76 & 70.69 \\
\bottomrule
\end{tabular}
\caption{Code-mixed sentiment classification on \textsc{Theedhum Nandrum}. The first three columns replace the input with its transliteration; the last concatenates it (augmentation). The two setups disagree in sign on Malayalam macro-F1---baseline wins under replacement, phonology-aware wins under augmentation---and every gap is under $1.4$ points on a test set of roughly $4.5$k examples. This indicates no reliable downstream effect in either direction rather than evidence for either scheme.}
\label{tab:downstream}
\end{table}